% ===== PRÉAMBULE CANONIQUE — à coller VERBATIM en tête de CHAQUE .tex =====
\documentclass[11pt,a4paper]{article}
\usepackage[utf8]{inputenc}\usepackage[T1]{fontenc}\usepackage{lmodern}
\usepackage[french]{babel}
\usepackage{amsmath,amssymb,mathtools}\usepackage{icomma}
\usepackage[version=4]{mhchem}          % \ce{...} : équations & formules
\usepackage{siunitx}\sisetup{locale=FR} % \qty{}{}, \si{}, \ang{}
\usepackage{chemfig}                    % \chemfig{...} : structures développées
\usepackage{xcolor}\usepackage{colortbl}
\usepackage{tikz}\usepackage{pgfplots}\pgfplotsset{compat=1.18}
\usepackage[most]{tcolorbox}\usepackage{enumitem}\usepackage{array,booktabs}
\usepackage[a4paper,margin=2.2cm]{geometry}
\usetikzlibrary{arrows.meta,calc,angles,quotes,positioning,patterns,backgrounds,shapes.geometric}
\definecolor{primary}{RGB}{222,49,99}\definecolor{primarydark}{RGB}{150,22,55}
\definecolor{defcol}{RGB}{0,114,178}\definecolor{methcol}{RGB}{52,92,148}
\definecolor{excol}{RGB}{0,135,90}\definecolor{attncol}{RGB}{200,40,55}
\tcbset{boxbase/.style={enhanced,breakable,boxrule=0.9pt,arc=2.5pt,
  left=3mm,right=3mm,top=2.3mm,bottom=2.3mm,fonttitle=\bfseries,coltitle=white}}
% --- boîtes TUTORIEL (noms EXACTS, ne pas renommer) ---
\newtcolorbox{cle}[1][]{boxbase,colback=defcol!6,colframe=defcol,
  title={\textsc{À retenir}\ifx\relax#1\relax\relax\else\textmd{ : \itshape #1}\fi}}
\newtcolorbox{methode}[1][]{boxbase,colback=methcol!7,colframe=methcol,sharp corners,
  title={\textsc{Méthode}\ifx\relax#1\relax\relax\else\textmd{ --- \itshape #1}\fi}}
\newtcolorbox{exemplebox}[1][]{boxbase,colback=excol!5,colframe=excol,
  title={\textsc{Exemple}\ifx\relax#1\relax\relax\else\textmd{ : \itshape #1}\fi}}
\newtcolorbox{piege}[1][]{boxbase,colback=attncol!6,colframe=attncol,
  title={\textsc{Piège}\ifx\relax#1\relax\relax\else\textmd{ : \itshape #1}\fi}}
\newtcolorbox{remarque}[1][]{boxbase,colback=black!4,colframe=black!45,
  title={\textsc{Remarque}\ifx\relax#1\relax\relax\else\textmd{ : \itshape #1}\fi}}
% --- boîtes EXERCICE / CORRIGÉ (noms EXACTS, auto-comptées) ---
\newtcolorbox[auto counter]{exo}[1][]{boxbase,colback=primary!4,colframe=primary,
  title={\textsc{Exercice}~\thetcbcounter\ifx\relax#1\relax\relax\else\textmd{ --- \itshape #1}\fi}}
\newtcolorbox[auto counter]{corrige}[1][]{boxbase,colback=excol!4,colframe=excol,
  title={\textsc{Corrigé}~\thetcbcounter\ifx\relax#1\relax\relax\else\textmd{ --- \itshape #1}\fi}}
\newcommand{\R}{\mathbb{R}}\newcommand{\N}{\mathbb{N}}\newcommand{\Z}{\mathbb{Z}}\newcommand{\ds}{\displaystyle}
% (un \usepackage{fancyhdr}+en-tête est possible ; il est ignoré côté web)
% =========================================================================
\begin{document}

\begin{corrige}[le piège de l'ion phosphite]
\begin{enumerate}[label=\alph*)]
  \item ion phosphite : \ce{HPO3^2-}.
  \item ion phosphate : \ce{PO4^3-}.
  \item parce que l'acide phosphoreux \ce{H3PO3} ne libère que \emph{deux} \ce{H+} (il est diprotique) : en
        perdant ses deux \ce{H} acides, il donne \ce{HPO3^2-}, où subsiste le \ce{H} non dissociable lié au
        phosphore. La formule « \ce{PO3^3-} » supposerait à tort trois \ce{H} acides.
\end{enumerate}
\end{corrige}

\begin{corrige}[la cascade de l'acide phosphorique]
On préfixe « (di)(mono)hydrogéno » selon le nombre de \ce{H} restants.
\begin{enumerate}[label=\alph*)]
  \item \ce{H2PO4-} : dihydrogénophosphate.
  \item \ce{HPO4^2-} : monohydrogénophosphate (ou hydrogénophosphate).
  \item \ce{PO4^3-} : phosphate.
\end{enumerate}
\end{corrige}

\begin{corrige}[de l'oxacide à son anhydride]
On soustrait \ce{H2O} et l'on regroupe les atomes restants.
\begin{enumerate}[label=\alph*)]
  \item $\ce{H2SO4} - \ce{H2O} = \ce{SO3}$ (soufre $+\mathrm{VI}$ des deux côtés).
  \item $\ce{2 HNO3} - \ce{H2O} = \ce{N2O5}$ (azote $+\mathrm{V}$).
  \item $\ce{2 HClO4} - \ce{H2O} = \ce{Cl2O7}$ (chlore $+\mathrm{VII}$).
\end{enumerate}
\end{corrige}

\begin{corrige}[molécule (acide) ou ion ?]
Le signe de charge tranche : sans charge $\Rightarrow$ molécule ; avec charge $\Rightarrow$ ion.
\begin{enumerate}[label=\alph*)]
  \item \ce{H2SO3} : molécule neutre $\Rightarrow$ acide sulfureux (ou sulfite d'hydrogène).
  \item \ce{SO3^2-} : ion $\Rightarrow$ ion sulfite.
  \item \ce{HSO3-} : ion $\Rightarrow$ ion hydrogénosulfite.
\end{enumerate}
\end{corrige}

\begin{corrige}[synthèse : chromique, dichromique et n.o.]
\begin{enumerate}[label=\alph*)]
  \item acide chrom\emph{ique} $\to$ anion chrom\emph{ate} \ce{CrO4^2-}.
  \item acide dichrom\emph{ique} $\to$ anion dichrom\emph{ate} \ce{Cr2O7^2-}.
  \item \ce{CrO4^2-} : $x+4(-2)=-2 \Rightarrow x=+\mathrm{VI}$ ;\quad
        \ce{Cr2O7^2-} : $2x+7(-2)=-2 \Rightarrow x=+\mathrm{VI}$. Le chrome est $+\mathrm{VI}$ dans les deux.
\end{enumerate}
\end{corrige}

\end{document}
